% ===================== THE SHAPES OF GALAXIES =====================
% Twenty-one galaxy types, two to a page in stacked rows (top/bottom),
% three on the closing page. Ten pages over the starfield plate.

One more step out, and the object of study is no longer a star or a system but
an entire island of them. Twenty-one kinds are catalogued below: the shapes
Hubble sorted a century ago, the wrecks and rings only collisions explain, the
extremes of size and darkness, the ones whose cores are eating, and the ones
burning their gas too fast to last.

\galaxyfieldon
\galaxywide

\galaxyhead{The Shapes of Galaxies}

% ---------- 1. THE STANDARD HUBBLE SEQUENCE ----------

\galaxycard{uf_gal_spiral}{Spiral Galaxies (S)}{Hubble sequence; Messier 74,
32 million light years}{A rotating disc of gas, dust and young stars wound into
curved arms around a central bulge. The arms are not solid objects but density
waves: material piles up as it drifts through them, and the pile-up is where new
stars light. The Milky Way is one of these, which is why the class matters more
than the rest --- everything we know about living inside a galaxy is measured
from inside this shape.}{NASA, ESA, and the Hubble Heritage (STScI/AURA)-ESA/Hubble Collaboration. Acknowledgment: R. Chandar (University of Toledo) and J. Miller (University of Michigan). Public domain.}{https://commons.wikimedia.org/wiki/File:Messier_74_by_HST.jpg}

\galaxycard{uf_gal_barred}{Barred Spiral Galaxies (SB)}{Hubble sequence; NGC
1300, 61 million light years}{A spiral whose arms spring not from a round
nucleus but from the ends of a straight bar of stars driven through the centre.
The bar is a gravitational instability of the disc itself, and it works as a
funnel: it drags gas inward to feed the core. Our own galaxy has one, roughly
27,\!000 light years long, which we could not see from inside until infrared
surveys traced the star counts through the dust.}{NASA, ESA, and The Hubble Heritage Team (STScI/AURA). Public domain.}{https://commons.wikimedia.org/wiki/File:Hubble2005-01-barred-spiral-galaxy-NGC1300.jpg}

\galaxybreak
\galaxycard{uf_gal_elliptical}{Elliptical Galaxies (E)}{Hubble sequence;
Messier 87, 53 million light years}{A smooth featureless swarm, spherical to
football-shaped, made almost entirely of old red stars with very little gas left
to make new ones. Ellipticals are what you get after enough mergers: the orderly
rotation of the parent discs is scrambled into random orbits. M87 is the giant
at the heart of the Virgo cluster and the host of the first black hole ever
photographed.}{en:NASA, en:STScI, en:WikiSky. Public domain.}{https://commons.wikimedia.org/wiki/File:Messier_87_Hubble_WikiSky.jpg}

\galaxycard{uf_gal_lenticular}{Lenticular Galaxies (S0/SB0)}{Hubble sequence;
NGC 5866, 44 million light years}{The transitional case: a galaxy that still has
the bulge and the disc of a spiral but has lost or used up the gas, so there are
no arms left to trace. Edge-on, as here, the dust lane is a knife-line across the
centre. Lenticulars sit at the hinge of the Hubble tuning fork and are the best
evidence that his sequence is not an evolutionary ladder but a family of
end-states.}{NASA, ESA, and The Hubble Heritage Team (STScI/AURA). Public domain.}{https://commons.wikimedia.org/wiki/File:Ngc5866_hst_big.jpg}

\galaxybreak
\galaxycard{uf_gal_irregular}{Irregular Galaxies (Irr)}{Hubble sequence; Large
Magellanic Cloud, 163,\!000 light years}{No disc, no bulge, no symmetry --- just
a chaotic mass of gas and star-forming knots, usually because a larger neighbour
has pulled the shape apart. The Large Magellanic Cloud is the nearest example and
is being stretched by the Milky Way as it falls in. Irregulars are common in the
early universe, which suggests order in galaxies is something that accumulates
rather than something they start with.}{ESO, IAU and Sky \& Telescope. CC BY 4.0.}{https://commons.wikimedia.org/wiki/File:Large_Magellanic_Cloud_(eso1914c).jpg}

% ---------- 2. EXOTIC AND RARE STRUCTURAL TYPES ----------

\galaxycard{uf_gal_ring}{Ring Galaxies}{Exotic structure; Hoag's Object, 600
million light years}{A bright dense core, then a gap, then an isolated ring of
brilliant blue young stars. The standard explanation is a bullseye collision: a
smaller galaxy punches through the centre of a disc and the shock drives a
circular wave of star formation outward. Hoag's Object is the textbook case and
still argued over, because no obvious intruder galaxy has ever been found nearby
to have done the punching.}{NASA and The Hubble Heritage Team (STScI/AURA); Acknowledgment: Ray A. Lucas (STScI/AURA). Public domain.}{https://commons.wikimedia.org/wiki/File:Hoag\%27s_object.jpg}

\galaxybreak
\galaxycard{uf_gal_peculiar}{Peculiar Galaxies}{Exotic structure; the Antennae,
NGC 4038/4039, 60 million light years}{Two galaxies in the act of merging, with
tidal tails thrown hundreds of thousands of light years out and the gas between
them collapsing into new clusters. Stars almost never hit each other in a
collision like this --- the space between them is far too large --- but the gas
clouds do, which is why a wreck this violent looks bright rather than dark. This
is the Milky Way and Andromeda in about four billion years.}{john.purvis. CC BY 2.0.}{https://commons.wikimedia.org/wiki/File:HLA_data_-_The_Antennae_Galaxies_-_NGC4038_4039_(Arp_244)_in_Corvus_(39931666833).jpg}

\galaxycard{uf_gal_polarring}{Polar-Ring Galaxies}{Exotic structure; NGC 4650A,
130 million light years}{A rare double system: an inner disc rotating one way and
an outer ring of gas and stars rotating over its poles, perpendicular to it. The
ring is captured material --- either a shredded companion or gas pulled from one
--- and because the two components orbit in different planes, the shape is a
direct measurement of the dark matter halo's geometry. Very few clean examples are
known.}{The Hubble Heritage Team (AURA/STScI/NASA). Public domain.}{https://commons.wikimedia.org/wiki/File:NGC_4650A_I_HST2002.jpg}

\galaxybreak
\galaxycard{uf_gal_shell}{Shell Galaxies}{Exotic structure; NGC 3923, 68 million
light years}{An elliptical wrapped in faint concentric arcs of stars, nested one
inside the next like the layers of an onion. The shells are the debris of a
smaller galaxy that fell in and was torn apart, its stars settling into ordered
ripples at the turning points of their orbits. Counting the shells dates the meal.
NGC 3923 has more than twenty and the outer ones are enormous.}{NASA/ESA Hubble Space Telescope. Public domain.}{https://commons.wikimedia.org/wiki/File:NGC_3923_Elliptical_Shell_Galaxy.jpg}

% ---------- 3. SIZE AND DENSITY VARIETIES ----------

\galaxycard{uf_gal_giant}{Giant Galaxies}{Size class; IC 1101, roughly 1 billion
light years}{The largest structures made of stars: supergiant ellipticals sitting
at the centres of galaxy clusters, grown by swallowing everything that falls in.
IC 1101 is the usual claimant, with a diffuse halo measured in the millions of
light years and an estimated star count of order a hundred trillion. Its
boundaries are genuinely disputed, because a galaxy that faint at the edge has no
agreed edge.}{NASA/ESA/Hubble Space Telescope. Public domain.}{https://commons.wikimedia.org/wiki/File:IC_1101_in_Abell_2029_(hst_06228_03_wfpc2_f702w_pc).jpg}

\galaxybreak
\galaxycard{uf_gal_dwarf}{Dwarf Galaxies}{Size class; NGC 1569, 11 million light
years}{A few thousand to a few billion stars, usually in orbit around a larger
parent. They are by far the most numerous kind of galaxy in the universe and the
hardest to find, because low mass means low surface brightness. The Milky Way has
dozens of known satellites and almost certainly many more undetected. Dwarfs are
also the building blocks: bigger galaxies are made by eating these.}{ESA, NASA and P. Anders (G\"ottingen University Galaxy Evolution Group, Germany). Public domain.}{https://commons.wikimedia.org/wiki/File:Supernova_Bonanza_in_Nearby_Galaxy_NGC_1569_(2004-06-1455).jpg}

\galaxycard{uf_gal_udg}{Ultra-Diffuse Galaxies (UDGs)}{Size class; NGC 1052-DF4,
45 million light years}{Ghosts: as wide across as the Milky Way but with something
like one per cent of its stars, so faint that most were missed until wide-field
surveys in the last decade. Some appear to be held together by overwhelming dark
matter; a few, including the DF2 and DF4 systems imaged here, appear to have
almost none, which is a real and unresolved problem for the standard model of
galaxy formation.}{European Space Agency. CC BY 2.0.}{https://commons.wikimedia.org/wiki/File:New_Hubble_data_explains_missing_dark_matter_(50652188762).jpg}

\galaxybreak
\galaxycard{uf_gal_ucd}{Ultra-Compact Dwarfs (UCDs)}{Size class; M60-UCD1, 54
million light years}{The opposite extreme: millions of stars packed into a few
hundred light years, so dense that the night sky from a planet inside one would
hold a million naked-eye stars. M60-UCD1 also carries a supermassive black hole
amounting to a large fraction of its mass, which is the strongest evidence that
UCDs are not star clusters at all but the stripped nuclei of galaxies that lost
everything else.}{NASA, ESA, CXC, and J. Strader (Michigan State University). Public domain.}{https://commons.wikimedia.org/wiki/File:M60-UCD1_by_HST.jpg}

% ---------- 4. ACTIVE AND ENERGETIC VARIETIES (AGN) ----------

\galaxycard{uf_gal_quasar}{Quasars}{Active nucleus; 3C 273-class core, imaged in
nearby quasar host galaxies}{A galactic core whose central black hole is accreting
so hard that the infalling matter outshines every star in the host combined. The
name is a fossil of the 1960s, when they looked like stars with impossible
redshifts. Because the brightest are the most distant objects we can routinely
see, they double as backlights: their spectra record every cloud of gas the light
crossed on the way here.}{WFPC2: NASA and J. Bahcall (IAS), A. Martel (JHU), H. Ford (JHU), M. Clampin (STScI), G. Hartig (STScI), G. Illingworth (UCO/Lick Observatory), the ACS Science Team and ESA. Public domain.}{https://commons.wikimedia.org/wiki/File:Hubble_probes_the_heart_of_a_nearby_quasar_(opo0303b).jpg}

\galaxybreak
\galaxycard{uf_gal_blazar}{Blazars}{Active nucleus; artist's concept}{The same
engine as a quasar, but with one of its plasma jets pointed almost exactly at
Earth. Looking down the barrel compresses the timescales and amplifies the
brightness, so blazars flicker violently across gamma rays in hours. No telescope
can resolve the jet base of one, so the honest picture is a NASA illustration
rather than a photograph --- there is nothing to photograph but an unresolved
point.}{NASA/JPL-Caltech/GSFC. Public domain. Artist's concept, not a photograph.}{https://commons.wikimedia.org/wiki/File:Blazar_-_Artist_Concept.tif}

\galaxycard{uf_gal_seyfert}{Seyfert Galaxies (Types I and II)}{Active nucleus;
NGC 7742, 72 million light years}{An ordinary-looking spiral with a hyper-luminous
core pouring out infrared, ultraviolet and X-rays. Seyferts are the nearby,
lower-powered relatives of quasars, which makes them the class where the accretion
physics can actually be studied. Type I shows broad emission lines from gas close
in; Type II hides that region behind dust, and the difference is now read mostly
as viewing angle.}{IMAGE: AURA, STScI, NASA, Hubble Heritage Project (STScI, AURA). Public domain.}{https://commons.wikimedia.org/wiki/File:Black_Hole-Powered_Spiral_Galaxy_NGC_7742_(1998-28-696).jpg}

\galaxybreak
\galaxycard{uf_gal_radio}{Radio Galaxies}{Active nucleus; Centaurus A, 12 million
light years}{Galaxies that throw invisible plumes of radio-emitting plasma far
beyond their own starlight, sometimes for millions of light years. In visible light
Centaurus A is an elliptical cut by a dust lane; in submillimetre and X-ray, as
composited here, the jets erupting from its core dominate the frame. This is the
nearest example and the reason radio astronomy discovered that galaxies do things
optical astronomy cannot see.}{ESO/WFI (Optical); MPIfR/ESO/APEX/A. Weiss et al. (Submillimetre); NASA/CXC/CfA/R. Kraft et al. (X-ray). CC BY 4.0.}{https://commons.wikimedia.org/wiki/File:ESO_Centaurus_A_LABOCA.jpg}

\galaxycard{uf_gal_liner}{LINER Galaxies}{Active nucleus; NGC 4036, 70 million
light years}{Low-Ionization Nuclear Emission-line Regions: cores that are mildly,
unspectacularly active, showing emission lines from weakly ionised gas rather than
the blaze of a Seyfert. They are extremely common --- perhaps a third of nearby
galaxies qualify --- and the cause is still contested between a starved black hole
and ordinary hot old stars. A useful reminder that ``active'' is a spectrum, not a
switch.}{ESA/Hubble \& NASA. Acknowledgement: Judy Schmidt. CC BY 4.0.}{https://commons.wikimedia.org/wiki/File:Hazy_dust_in_Ursa_Major_NGC_4036.jpg}

% ---------- 5. STAR-FORMING ABNORMALITIES ----------

\galaxybreak
\galaxytrio
\galaxycard{uf_gal_starburst}{Starburst Galaxies}{Star formation; Messier 82, 12
million light years}{A galaxy converting gas into stars at a rate it cannot
possibly sustain, usually because a close pass with a neighbour compressed its
clouds all at once. M82 is making stars roughly ten times faster than the Milky Way
and venting the exhaust: the red filaments here are hydrogen being blown out of the
disc by the combined winds and supernovae of the burst itself.}{NASA, ESA, and The Hubble Heritage Team (STScI/AURA). Public domain.}{https://commons.wikimedia.org/wiki/File:M82_HST_ACS_2006-14-a-large_web.jpg}

\galaxycard{uf_gal_greenpea}{Green Pea Galaxies}{Star formation; illustration of
the class}{Small round compact galaxies that appear green in survey images because
intense doubly-ionised oxygen emission lands in the green filter. They were found
in 2007 by Galaxy Zoo volunteers, not by professionals, and they matter because
their furious bursts leak ultraviolet light --- making them the closest available
analogue to the galaxies that reionised the early universe. Every real Green Pea is
a featureless dot at this scale, so the card carries an illustration.}{geckzilla. CC BY 2.0. Illustration of the class, not a telescope image.}{https://commons.wikimedia.org/wiki/File:Green_Pea_Galaxy_Illustration_(29801410088).png}

\galaxycard{uf_gal_bcd}{Blue Compact Dwarfs (BCDs)}{Star formation; NGC 1705, 17
million light years}{Small, gas-rich, chemically primitive galaxies whose light is
dominated by clusters of hot young massive stars, which tints the whole system
blue. NGC 1705 holds a super star cluster near its centre and is driving a
galactic wind out of it. Because their gas has been through few generations of
stars, BCDs are used as local stand-ins for what the first galaxies were made
of.}{NASA, ESA, and The Hubble Heritage Team (STScI/AURA). Public domain.}{https://commons.wikimedia.org/wiki/File:NGC_1705.jpg}

\galaxynarrow
\galaxyfieldoff
